\section{Related Work}

\paragraph{Minecraft as a Test-bed for AI}
Minecraft has gained popularity as a benchmark for AI research due to its complex and dynamic environment, making it a rich test-bed for reinforcement learning and other AI methods (e.g., \citep{johnson2016malmo, guss2019minerl,fan2022minedojo, hafner2023dreamerv3,nottingham2023embodied, wang2023describe,malato2022behavioral,cai2023open}). We leverage the MineRL environment \citep{guss2019minerl} to research the creation of agents that can follow open-ended instructions in complex visual environments using only low-level actions (mouse and keyboard). We build \agentname{} on top of two recent foundation models. In order to align text and videos, we use MineCLIP \citep{fan2022minedojo}, a CLIP \citep{radford2021clip} model trained on paired web videos of Minecraft gameplay and associated captions. To train \agentname{}'s policy, we fine-tune VPT \citep{baker2022vpt}, a foundation model of Minecraft behavior that is pretrained on 70k hours of web videos of Minecraft along with estimated mouse and keyboard~actions. Several prior works \citep{volum-etal-2022-craft, wang2023describe} have explored the use of LLMs in creating instructable Minecraft agents. These works typically use LLMs to make high-level plans that are then executed by lower-level RL \citep{nottingham2023embodied, wang2023describe} or scripted \citep{PrismarineJS2023} policies. Since \agentname{} is a far more flexible low-level policy, the combination of \agentname{} with LLMs is a promising direction for future work. \citet{fan2022minedojo} introduced an agent trained using RL with MineCLIP as a shaping reward on 12 different tasks and conditioned on MineCLIP-embedded text-prompts. However, this agent failed to generalize beyond the original set of tasks without further RL finetuning using the MineCLIP reward function. \citet{cai2023open} proposed a Goal-Sensitive Backbone architecture for goal-conditioned control in Minecraft which is trained on a fixed set of goals, while \agentname{} learns goal-reaching behavior from a large dataset in a self-supervised way without training on an explicit set of tasks.


\paragraph{Foundation Models for Sequential Decision-Making}
Foundation models which are pretrained on vast amounts of data and then finetuned for specific tasks have recently shown great promise in a variety of domains including language \citep{brown2020language, palm, llama}, vision \citep{ramesh2022dalle2, caron2021emerging, radford2021clip}, and robotics \citep{brohan2022rt, shridhar2022cliport, jiang2022vima, nair2022r3m, xiao2022robotic}.
GATO \citep{reed2022gato} and RT-1 \citep{brohan2022rt} have demonstrated the potential of training transformers to perform both simulated and real-world robotic tasks. With the exception of \citet{kumar2023offline}, which uses Q-learning, the vast majority of cases \citep{lee2022multi, brohan2022rt, reed2022gato} where deep learning has been scaled to large, multitask offline-RL datasets have used supervised RL. Supervised RL (e.g., \citep{paster2020planning, ghosh2021gcsl, chen2021decision}) works by framing the sequential decision-making problem as a prediction problem, where the model is trained to predict the next action conditioned on some future outcome. While these approaches are simple and scale well with large amounts of compute and data, more work is needed to understand the trade-offs between supervised RL and Q-learning or policy gradient-based methods \citep{paster2022you, paster2022return, brandfonbrener2022does, udrl_stochastic}. Recent works explore the use of hindsight relabeling \citep{hindsight} using vision-language models \citep{radford2021clip, alayrac2022flamingo} to produce natural language relabeling instructions. DIAL \citep{xiao2022robotic} finetunes CLIP \citep{radford2021clip} on human-labeled trajectories, which is then used to select a hindsight instruction from a candidate set. \citet{sumers2023distilling} uses Flamingo \citep{alayrac2022flamingo} zero-shot for hindsight relabeling by framing it as a visual-question answering (VQA) task. In contrast, \agentname{} relabels goals using future trajectory segment embeddings given by the MineCLIP \citep{fan2022minedojo} visual embedding.



\paragraph{Text-Conditioned Generative Models}
There has been a recent explosion of interest in text-to-X models, including text-to-image (e.g., \citep{ramesh2022dalle2, imagen, rombach2022ldm}), text-to-3D (e.g., \citep{shapee, lin2023magic3d}), and even text-to-music (e.g., \citep{musiclm}). These models are typically either autoregressive transformers modeling sequences of discrete tokens \citep{transformer, brown2020language} or diffusion models \citep{diffusion}. Most related to our work is unCLIP, the method used for DALL•E 2 \citep{ramesh2022dalle2}. unCLIP works by training a generative diffusion model to sample images from CLIP \citep{radford2021clip} embeddings of those images. By combining this model with a prior that translates text to visual CLIP embeddings, unCLIP can produce photorealistic images for arbitrary text prompts. unCLIP and many other diffusion-based approaches utilize a technique called classifier-free guidance \citep{cfg}, which lets the model trade-off between mode-coverage and sample fidelity post-training. We utilize the basic procedure of unCLIP and classifier-free guidance for training \agentname{}.